\chapter{级联二选一数据选择器}

\section{第一步：解决什么问题}

单个2选1 MUX只能从2路数据中选择1路。如果要从更多路数据中选择，怎么办？

\textbf{方案}：用多个2选1 MUX级联，构建更大的MUX。

\begin{keypoint}[级联MUX的本质]
用最基础的2选1 MUX作为构建单元，通过\textbf{分层选择}的方式扩展选择能力。

2选1 $\to$ 4选1 $\to$ 8选1 $\to$ 16选1 $\to$ ...
\end{keypoint}

\section{第二步：用2选1 MUX构建4选1 MUX}

先看2选1 MUX的基本符号：$Y = \overline{S} \cdot D_0 + S \cdot D_1$

\textbf{用3个2选1 MUX构建4选1 MUX}：

\begin{center}
\begin{tikzpicture}
  % First layer: two 2:1 MUXes
  \node[draw, minimum width=2cm, minimum height=0.8cm] (mux0) at (0,2) {MUX 2:1};
  \node[draw, minimum width=2cm, minimum height=0.8cm] (mux1) at (0,0) {MUX 2:1};
  % Second layer: one 2:1 MUX
  \node[draw, minimum width=2cm, minimum height=0.8cm] (mux2) at (4,1) {MUX 2:1};

  % Inputs to first layer
  \node (D0) at (-3,2.5) {$D_0$}; \draw[->] (D0) -- (mux0);
  \node (D1) at (-3,1.5) {$D_1$}; \draw[->] (D1) -- (mux0);
  \node (D2) at (-3,0.5) {$D_2$}; \draw[->] (D2) -- (mux1);
  \node (D3) at (-3,-0.5) {$D_3$}; \draw[->] (D3) -- (mux1);

  % First layer outputs to second layer
  \draw[->] (mux0.east) -- ++(0.5,0) |- (mux2.west |- 2);
  \draw[->] (mux1.east) -- ++(0.5,0) |- (mux2.west |- 0);

  % Selection
  \node (S0) at (-2,-2) {$S_0$};
  \draw[->] (S0) -- (0,-1.5) node[above] {};

  \node (S1) at (4,-2) {$S_1$};

  % Output
  \node (Y) at (7,1) {$Y$};
  \draw[->] (mux2.east) -- (Y);
\end{tikzpicture}
\end{center}

\textbf{工作原理}：
\begin{itemize}
  \item $S_0$ 控制第一层两个MUX：分别从 $(D_0,D_1)$ 和 $(D_2,D_3)$ 中各选一个
  \item $S_1$ 控制第二层MUX：从第一层两个结果中选一个
  \item 结果：$(S_1,S_0)$ 决定了选择 $D_0, D_1, D_2, D_3$ 中的哪一个
\end{itemize}

\section{第三步：真值表}

\begin{center}
\begin{tabular}{|c|c||c|c|c|}
\hline
$S_1$ & $S_0$ & 第一层左MUX输出 & 第一层右MUX输出 & 最终输出 $Y$ \\
\hline
0 & 0 & $D_0$ & $D_2$ & $D_0$ \\
0 & 1 & $D_1$ & $D_3$ & $D_1$ \\
1 & 0 & $D_0$ & $D_2$ & $D_2$ \\
1 & 1 & $D_1$ & $D_3$ & $D_3$ \\
\hline
\end{tabular}
\end{center}

\section{第四步：逻辑推导}

级联MUX的逻辑表达式可以通过代入法推导：

第一层输出：
\begin{align}
  F_0 &= \overline{S_0} \cdot D_0 + S_0 \cdot D_1 \\
  F_1 &= \overline{S_0} \cdot D_2 + S_0 \cdot D_3
\end{align}

第二层输出：
\begin{align}
  Y &= \overline{S_1} \cdot F_0 + S_1 \cdot F_1 \\
    &= \overline{S_1} \cdot (\overline{S_0} \cdot D_0 + S_0 \cdot D_1)
     + S_1 \cdot (\overline{S_0} \cdot D_2 + S_0 \cdot D_3) \\
    &= \overline{S_1}\cdot\overline{S_0}\cdot D_0
     + \overline{S_1}\cdot S_0 \cdot D_1
     + S_1\cdot\overline{S_0} \cdot D_2
     + S_1\cdot S_0 \cdot D_3
\end{align}

\textbf{这正是标准4选1 MUX的逻辑表达式！} 证明级联等价于一个完整的$2^n$选1 MUX。

\section{第五步：结构化设计——通用级联规律}

\textbf{树形结构}：

\begin{itemize}
  \item \textbf{第一层}：$2^{n-1}$ 个2选1 MUX（用最低位 $S_0$ 选择）
  \item \textbf{第二层}：$2^{n-2}$ 个2选1 MUX（用 $S_1$ 选择）
  \item \textbf{...}
  \item \textbf{第n层}：1个2选1 MUX（用最高位 $S_{n-1}$ 选择）
\end{itemize}

总MUX数量：$2^n - 1$ 个2选1 MUX。

\section{第六步：为什么这样设计}

\begin{enumerate}
  \item \textbf{为什么2选1是基础构建块？}

  2选1 MUX是最简单的MUX（1位选择线），在FPGA中通常对应一个LUT（查找表）。任何更大的MUX都可以由2选1 MUX级联构建。

  \item \textbf{树形结构的深层意义？}

  每一层将数据组数减半：
  \begin{itemize}
    \item 输入：$2^n$ 路数据
    \item 第一层后：$2^{n-1}$ 路
    \item 第二层后：$2^{n-2}$ 路
    \item ...
    \item 第n层后：1路（最终输出）
  \end{itemize}

  这本质上是锦标赛淘汰赛的结构——每一轮淘汰一半。

  \item \textbf{延迟分析}：

  级联MUX的关键路径（从数据输入到输出）经过 $n$ 级2选1 MUX。每级延迟约1-2个门延迟。
\end{enumerate}

\section{第七步：考试常见变式}

\begin{enumerate}
  \item \textbf{变式1：用2选1 MUX构建8选1 MUX}

  需要 $2^3 - 1 = 7$ 个2选1 MUX。三层结构：4+2+1。

  \item \textbf{变式2：用4选1 MUX构建16选1 MUX}

  5片4选1 MUX：4片第一层 + 1片第二层（作为最终选择）。

  \item \textbf{变式3：不对称级联}

  例如：用1片4选1 + 1片2选1构建6选1 MUX（实际数据输入不足$2^n$个时的处理方法）。

  \item \textbf{变式4：MUX树实现逻辑函数}

  将逻辑函数的变量从低位到高位接选择线，用MUX树实现多变量函数。这是MUX做通用逻辑的关键应用。
\end{enumerate}

\section{第八步：VHDL实现}

\begin{lstlisting}[caption={用2选1 MUX级联构建4选1 MUX（结构级）}]
library ieee;
use ieee.std_logic_1164.all;

entity mux2to1 is
  port (
    D0, D1 : in  std_logic;
    S      : in  std_logic;
    Y      : out std_logic
  );
end mux2to1;

architecture behavioral of mux2to1 is
begin
  Y <= D0 when S = '0' else D1;
end behavioral;

-- ========================================

library ieee;
use ieee.std_logic_1164.all;

entity mux4to1_cascade is
  port (
    D : in  std_logic_vector(3 downto 0);
    S : in  std_logic_vector(1 downto 0);
    Y : out std_logic
  );
end mux4to1_cascade;

architecture structural of mux4to1_cascade is
  signal F0, F1 : std_logic;  -- 第一层输出

  component mux2to1 is
    port (D0, D1 : in std_logic; S : in std_logic; Y : out std_logic);
  end component;
begin
  -- 第一层：S(0)选择
  mux_a: mux2to1 port map (D(0), D(1), S(0), F0);
  mux_b: mux2to1 port map (D(2), D(3), S(0), F1);
  -- 第二层：S(1)选择
  mux_c: mux2to1 port map (F0, F1, S(1), Y);
end structural;
\end{lstlisting}

\section{第九步：Testbench验证}

\begin{lstlisting}[caption={级联MUX Testbench}]
library ieee;
use ieee.std_logic_1164.all;

entity mux4to1_cascade_tb is
end mux4to1_cascade_tb;

architecture test of mux4to1_cascade_tb is
  signal D : std_logic_vector(3 downto 0);
  signal S : std_logic_vector(1 downto 0);
  signal Y : std_logic;

  component mux4to1_cascade is
    port (
      D : in  std_logic_vector(3 downto 0);
      S : in  std_logic_vector(1 downto 0);
      Y : out std_logic
    );
  end component;
begin
  uut: mux4to1_cascade port map (D, S, Y);

  stimulus: process
  begin
    D <= "1010";
    S <= "00"; wait for 10 ns;  -- Y=D0=0
    S <= "01"; wait for 10 ns;  -- Y=D1=1
    S <= "10"; wait for 10 ns;  -- Y=D2=0
    S <= "11"; wait for 10 ns;  -- Y=D3=1
    wait;
  end process;
end test;
\end{lstlisting}

\section{第十步：如何快速解题}

\begin{examtip}[级联MUX快速解题法]
\begin{enumerate}
  \item \textbf{$n$选1 MUX所需2选1 MUX数目} = $n-1$
  \item \textbf{树形结构}：每层将数据组数减半
  \item \textbf{选择位分配}：最低位控制第一层，最高位控制最后一层
  \item \textbf{关键路径} = $n$级2选1 MUX延迟（$n = \log_2(N)$）
\end{enumerate}
\end{examtip}

\textbf{秒杀：用2选1 MUX构建任意MUX的公式}

$2^n$选1 MUX 需要 $(2^n - 1)$ 个2选1 MUX，分 $n$ 层。第 $k$ 层（从1开始）有 $2^{n-k}$ 个MUX，由 $S_{k-1}$ 控制。

\section{变式详解}
\chapter{级联MUX——4大变式详解}

\section{变式1：用2选1 MUX构建8选1 MUX}

\begin{detailbox}[完整教学]
\textbf{【树形结构】}3层：第一层4个→第二层2个→第三层1个=7个2选1 MUX。

\textbf{【选择位分配】}
\begin{itemize}
  \item S0→第一层(4个MUX)：每两组数据中选一个
  \item S1→第二层(2个MUX)：第一层结果中两两选择
  \item S2→第三层(1个MUX)：第二层两个结果中最终选择
\end{itemize}

\textbf{【通用公式】}$2^n$选1 MUX需要$2^n-1$个2选1 MUX，分$n$层。

\textbf{【VHDL——generate结构级】}
\begin{lstlisting}[caption={8选1 MUX（2选1树形级联）}]
library ieee;
use ieee.std_logic_1164.all;

entity mux8to1_cascade is
  port (
    D : in  std_logic_vector(7 downto 0);
    S : in  std_logic_vector(2 downto 0);
    Y : out std_logic
  );
end mux8to1_cascade;

architecture behavioral of mux8to1_cascade is
  -- 递归描述流水线MUX树
  signal layer0 : std_logic_vector(3 downto 0);
  signal layer1 : std_logic_vector(1 downto 0);
begin
  gen0: for i in 0 to 3 generate
    layer0(i) <= D(2*i) when S(0)='0' else D(2*i+1);
  end generate;
  gen1: for i in 0 to 1 generate
    layer1(i) <= layer0(2*i) when S(1)='0' else layer0(2*i+1);
  end generate;
  Y <= layer1(0) when S(2)='0' else layer1(1);
end behavioral;
\end{lstlisting}
\end{detailbox}


\section{变式2：用4选1 MUX构建16选1 MUX}

\begin{detailbox}[完整教学]
\textbf{【结构】}第一层4片4选1 MUX，第二层1片4选1 MUX。共5片。

S1S0→第一层(4片4选1 MUX)。S3S2→第二层(1片4选1 MUX选择第一层输出之一)。
\end{detailbox}


\section{变式3：不对称级联（构建6选1 MUX）}

\begin{detailbox}[完整教学]
\textbf{【问题】}需要6选1，但MUX只有$2^n$选1。用8选1实现6选1（两个输入悬空或接地）。

\textbf{【VHDL——4选1+2选1】}
\begin{lstlisting}[caption={6选1 MUX（4选1+2选1级联）}]
library ieee;
use ieee.std_logic_1164.all;

entity mux6to1_asymmetric is
  port (
    D : in  std_logic_vector(5 downto 0);
    S : in  std_logic_vector(2 downto 0);
    Y : out std_logic
  );
end mux6to1_asymmetric;

architecture structural of mux6to1_asymmetric is
  signal mux4_out : std_logic;

  component mux4to1 is
    port (
      D : in  std_logic_vector(3 downto 0);
      S : in  std_logic_vector(1 downto 0);
      Y : out std_logic
    );
  end component;
begin
  -- 第一层：4选1处理D0-D3
  mux4: mux4to1 port map(D=>D(3 downto 0), S=>S(1 downto 0), Y=>mux4_out);
  -- 顶层：2选1在mux4_out和D4/D5间选择
  Y <= mux4_out when S(2)='0' else
       D(4) when S(1 downto 0)="00" else D(5);
  -- 注意：S2=1时S1S0只需表示2种状态即可
end structural;
\end{lstlisting}

\textbf{【简化方案】}直接用8选1 MUX，D6和D7接地(固定为0)。
\end{detailbox}


\section{变式4：MUX树实现逻辑函数}

\begin{detailbox}[完整教学]
\textbf{【本质】}MUX=通用逻辑实现器。MUX树=多层选择=分而治之。

\textbf{【Shannon展开定理】}$F(x_1,x_2,\dots,x_n)= \overline{x_1}\cdot F(0,x_2,\dots,x_n) + x_1\cdot F(1,x_2,\dots,x_n)$

这正是2选1 MUX的表达式！递归应用=构建MUX树。

\textbf{【结论】}任何组合逻辑函数都可用2选1 MUX树实现——这其实就是FPGA中LUT的工作原理。
\end{detailbox}
